\documentclass[a4paper]{ctexart}

\usepackage{amsmath, amssymb, amsthm}

\newtheorem{AI}{AI~-}

\usepackage{geometry}
\geometry{a4paper, left=2.5cm, right=2cm, top=2.5cm, bottom=2cm}

\usepackage[colorlinks=true, linkcolor=blue, urlcolor=blue, citecolor=blue]{hyperref}
\usepackage{graphicx}
\graphicspath{{imgs/}{images/}}

\usepackage{float}
\usepackage{tikz}
\usetikzlibrary{mindmap,shadows}


\newcommand{\airef}[1]{\hyperref[#1]{AI~-~\ref*{#1}}}

\title{不确定性的数学基础}
\author{陈钦}
\date{\today}


\begin{document}
\maketitle

\section{测度论与概率论}

\begin{AI}
    \label{AI:math-mathematization}
    “一门学问，只有当它成功地被数学化时，才算真正开始了对自身的研究。”
\end{AI}

要描述不确定性，其基础就是数理统计、概率论、随机过程等工具。
这些数学工具的基础则是集合论、拓扑学、测度论等。
\airef{AI:math-mathematization}的意思是说，只有当我们能够用数学工具来描述和分析一个领域时，
我们才能真正理解这个领域的本质。
对于不确定性，我们需要一个严密的数学框架来描述和分析它，
而测度论和概率论正是提供了这样一个框架。

基本的概率理论中，关于不确定性的基本数学形式为三元组 $(\Omega, \mathcal{F}, P)$ 。
其中 $\Omega$ 是样本空间，
$\mathcal{F}$ 是 $\Omega$ 上的 $\sigma$-代数，
$P$ 是定义在 $\mathcal{F}$ 上的概率测度。
在这个三元组中，
$\Omega$ 描述了所有可能的结果，
$\mathcal{F}$ 描述了事件的集合，如\eqref{eq:sigma-algebra}所示，
而 $P$ 则为每个事件分配了一个概率值，如\eqref{eq:probability-measure}所示。
因此，概率论提供了一个严密的框架来描述和分析不确定性。
在这个框架中，我们可以定义随机变量、概率分布、期望值、方差等概念，这些都是理解和量化不确定性的关键工具。

\begin{equation}
    \mathcal{F} = \{
    \emptyset,  A_1, A_2, \ldots,  A_n, \Omega
    \}
    \label{eq:sigma-algebra}
\end{equation}

\begin{equation}
    P: \mathcal{F} \to [0, 1] \subset \mathbb{R}, \quad P(\Omega) = 1, \quad P(\emptyset) = 0
    \label{eq:probability-measure}
\end{equation}

对古典概率中经典的投硬币问题，样本空间 $\Omega$ 可以定义为 $\{H, T\}$，其中 $H$ 表示正面，$T$ 表示反面。
事件集合 $\mathcal{F}$ 可以定义为 $\{\emptyset, \{H\}, \{T\}, \Omega\}$，
其中 $\emptyset$ 表示不发生任何事件，$\{H\}$ 表示投出正面，$\{T\}$ 表示投出反面，
$\Omega$ 表示投出正面或反面。
概率测度 $P$ 可以定义为 $P(\emptyset) = 0$，$P(\{H\}) = 0.5$，$P(\{T\}) = 0.5$，$P(\Omega) = 1$。
显然，有

\begin{equation}
    P(\Omega) \equiv P(\{H, T\}) = P(\{H\}) + P(\{T\}) = 0.5 + 0.5 = 1
\end{equation}

这对于所有的独立事件都成立。

\begin{equation}
    P\left(\bigcup_{i=1}^{n} A_i\right) = \sum_{i=1}^{n} P(A_i),
    A_i \cap A_j = \emptyset, \forall i \neq j, i, j = 1, 2, \ldots, n
\end{equation}

这里有一个必须注意的概念，就是$\mathcal{F}$中包含的是事件，而测度$P$是定义在事件上的函数。
因此，$P$的定义域是$\mathcal{F}$，而不是$\Omega$。
虽然$\Omega$是所有可能结果的集合，但我们并不直接对$\Omega$中的元素赋予概率，
而是对$\mathcal{F}$中的事件赋予概率。

对于硬币的例子，我们的测度函数$P(\emptyset) = 0$，$P(\{H\}) = 0.5$，$P(\{T\}) = 0.5$，$P(\Omega) = 1$，
而不是，$P(H) = 0.5$，$P(T) = 0.5$。
这样才能更加合适的表达事件：国徽或者菊花，也就是$\{H, T\}$这个事件及其概率。

我们用一个实数坐标系上一个区间的均匀分布的例子来进一步说明事件与事件结果的区别。
例如，我们随机选择一个点来切割一个长度为1的线段，将其分成两部分。
在这个例子中，样本空间 $\Omega$ 可以定义为 $[0, 1]$，表示线段上所有可能的切割点。

事件集合 $\mathcal{F}$ 的定义非常复杂。
例如，可以定义$[0, 0.5]$ 和 $(0.5, 1]$分别表示切割点在前半段和后半段的事件。
概率测度 $P$ 可以定义为 $P\left([0, 0.5]\right) = P(0 \le x \le 0.5) = 0.5$
和 $P\left((0.5,1]\right) = P(0.5 < x \le 1) = 0.5$，表示切割点在前半段和后半段的概率。
再比如，我们切割点在 $[0.25, 0.75]$ 这个事件的概率为 $P\left([0.25, 0.75]\right) = P(0.25 \le x \le 0.75) = 0.5$。
而$\Omega$ 中任何一个单个点（事件$x=0.25$用$\{0.25\}$表示）的概率都是 $P\left(\{0.25\}\right) = P(x = 0.25) = 0$。
这个连续随机事件的例子很好地说明了上面三元组的定义和事件与事件结果的区别。

通过这两个例子，我们就可以回过头来理解$\mathcal{F}$ 是 $\Omega$ 上的 $\sigma$-代数的意义了。

\begin{AI}
    $\sigma$-代数可以理解为“在概率空间中允许被讨论和赋予概率的一类事件集合”。

    更具体地说，设 $\mathcal{F}$ 是样本空间 $\Omega$ 的若干子集构成的集合，若它满足以下三条性质，就称 $\mathcal{F}$ 是一个 $\sigma$-代数：
    \begin{enumerate}
        \item $\Omega \in \mathcal{F}$（全集事件可测）；
        \item 若 $A \in \mathcal{F}$，则其补集 $A^c \in \mathcal{F}$（对补运算封闭）；
        \item 若 $A_1, A_2, \dots \in \mathcal{F}$，则 $\bigcup_{n=1}^{\infty} A_n \in \mathcal{F}$（对可数并封闭）。
    \end{enumerate}

    由这些性质还能推出：$\emptyset \in \mathcal{F}$，并且对可数交也封闭。直观上，$\sigma$-代数规定了“哪些事件是可测的”，从而保证概率测度 $P$ 可以在这些事件上被一致、无矛盾地定义。
\end{AI}

概率空间的概念是概率论的基础，它为我们提供了一个严密的框架来描述和分析不确定性。
通过定义样本空间、事件集合和概率测度，我们可以系统地研究随机现象，计算事件的概率，并分析随机变量的分布和统计特征。
这就自然地进一步引出随机变量、概率分布和概率密度等概念了。

\section{随机变量、概率分布与概率密度}

随机变量是定义在概率空间上的一个函数，它将样本空间中的每个结果映射到一个实数。
随机变量可以是离散的，也可以是连续的。

当一个量在$\Omega$上取值时，我们称它为一个随机变量。
随机变量可以是离散的，也可以是连续的。

当$\Omega \subseteq \mathbb{R}$时，我们通常使用概率分布函数CDF（\airef{AI:cdf}）来描述随机变量的分布。
概率分布函数 $F(x)$ 满足以下条件：

\begin{equation}
    F(x) = P(X \leq x), \quad \forall x \in \mathbb{R}
\end{equation}

\begin{AI}\label{AI:cdf}
    概率分布函数（CDF）是描述随机变量分布的函数。对于一个随机变量 $X$，其概率分布函数 $F(x)$ 定义为：

    \begin{equation}
        F(x) = P(X \leq x), \quad \forall x \in \mathbb{R}
    \end{equation}

    概率分布函数具有以下性质：
    \begin{enumerate}
        \item 非递减性：$F(x)$ 是一个非递减函数，即如果 $x_1 < x_2$，则 $F(x_1) \leq F(x_2)$。
        \item 极限性质：$\lim_{x \to -\infty} F(x) = 0$ 和 $\lim_{x \to \infty} F(x) = 1$。
              这表示当 $x$ 趋近于负无穷时，CDF 的值趋近于0；当 $x$ 趋近于正无穷时，CDF 的值趋近于1。
        \item 右连续性：$F(x)$ 是一个右连续函数，即对于任意 $x$，$\lim_{t \to x^+} F(t) = F(x)$。
    \end{enumerate}
\end{AI}

概率密度函数PDF（\airef{AI:pdf}）用来描述连续随机变量的分布。
概率密度函数 $f(x)$ 满足以下条件：

\begin{equation}
    f(x) \geq 0, \quad \forall x \in \mathbb{R}
\end{equation}

\begin{equation}
    \int_{-\infty}^{\infty} f(x) \, dx = 1
\end{equation}

\begin{AI}\label{AI:pdf}
    概率密度函数（PDF）是描述连续随机变量分布的函数。对于一个连续随机变量 $X$，其概率密度函数 $f(x)$ 满足以下条件：

    \begin{enumerate}
        \item 非负性：$f(x) \geq 0$ 对于所有 $x \in \mathbb{R}$ 都成立。
              这意味着概率密度函数的值不能为负数。
        \item 归一化：$\int_{-\infty}^{\infty} f(x) \, dx = 1$。
              这表示整个实数轴上的概率密度函数的积分必须等于1，确保了总概率为1。
    \end{enumerate}
\end{AI}

\begin{AI}\label{AI:pdf-cdf}
    对于任意区间 $[a, b]$，随机变量 $X$ 落在该区间内的概率可以通过积分计算得到：

    \begin{equation}
        P(a \leq X \leq b) = \int_{a}^{b} f(x) \, dx
    \end{equation}

    因此，概率密度函数提供了一种方式来量化连续随机变量在不同取值范围内的概率分布。
\end{AI}

\section{均值、方差与无偏估计}

随机变量的均值（期望值）和方差是描述其分布特征的重要统计量。
对于一个随机变量 $X$，其均值（期望值）定义为：

\begin{equation}
    \mu \equiv \mathbb{E}[X] = \int_{-\infty}^{\infty} x f(x) \, dx
\end{equation}

方差定义为：

\begin{equation}
    \sigma^2 \equiv \mathrm{Var}(X) = \mathbb{E}[{(X - \mu)}^2] = \int_{-\infty}^{\infty} {(x - \mu)}^2 f(x) \, dx
\end{equation}

无偏估计是指一个统计量的估计值等于其期望值。
如果一个变量$X$的$n$次样本均值$\bar{X}$满足$\mathbb{E}[\bar{X}] = \mu$，其中$\mu$是总体均值，那么$\bar{X}$就是$\mu$的无偏估计。

对于方差的无偏估计，我们通常使用样本方差 $S^2$，定义为：

\begin{equation}
    S^2 = \frac{1}{n-1} \sum_{i=1}^{n} {(X_i - \bar{X})}^2
    \label{eq:unbiased-variance-estimator}
\end{equation}

\begin{AI}
    与\eqref{eq:unbiased-variance-estimator}所给样本方差 $S^2$ 相对应的还有一个有偏估计 $\mathcal{T}^2$，定义为：
    \begin{equation}
        \mathcal{T}^2 = \frac{1}{n} \sum_{i=1}^{n} {(X_i - \bar{X})}^2
        \label{eq:biased-variance-estimator}
    \end{equation}

    样本方差 $S^2$ 是总体方差 $\sigma^2$ 的无偏估计，而 $\mathcal{T}^2$ 则是有偏估计。
    具体来说，$\mathbb{E}[S^2] = \sigma^2$，而 $\mathbb{E}[\mathcal{T}^2] = \frac{n-1}{n} \sigma^2$。

    为了完整推导，设 $X_1,\dots,X_n$ 独立同分布，且
    $\mathbb{E}[X_i]=\mu$，$\mathrm{Var}(X_i)=\sigma^2<\infty$。
    样本均值定义为
    \begin{equation}
        \bar{X}=\frac{1}{n}\sum_{i=1}^{n}X_i.
    \end{equation}

    先证明一个关键恒等式：
    \begin{equation}
        \sum_{i=1}^{n}{(X_i-\bar{X})}^2
        = \sum_{i=1}^{n}{(X_i-\mu)}^2 - n{(\bar{X}-\mu)}^2.
        \label{eq:variance-decomposition}
    \end{equation}

    证明如下。由
    \begin{equation}
        X_i-\bar{X}=(X_i-\mu)-(\bar{X}-\mu),
    \end{equation}
    两边平方并求和得
    \begin{equation}
        \begin{split}
            \sum_{i=1}^{n}{(X_i-\bar{X})}^2
             & = \sum_{i=1}^{n}{(X_i-\mu)}^2
            -2(\bar{X}-\mu)\sum_{i=1}^{n}(X_i-\mu)
            +\sum_{i=1}^{n}{(\bar{X}-\mu)}^2                       \\
             & = \sum_{i=1}^{n}{(X_i-\mu)}^2
            -2(\bar{X}-\mu)\cdot n(\bar{X}-\mu)
            +n{(\bar{X}-\mu)}^2                                    \\
             & = \sum_{i=1}^{n}{(X_i-\mu)}^2 - n{(\bar{X}-\mu)}^2.
        \end{split}
    \end{equation}
    因此\eqref{eq:variance-decomposition}成立。

    另外，独立同分布条件下
    \begin{equation}
        \begin{split}
            \mathrm{Var}(\bar{X})
             & = \mathrm{Var}\left(\frac{1}{n}\sum_{i=1}^{n}X_i\right)
            = \frac{1}{n^2}\mathrm{Var}\left(\sum_{i=1}^{n}X_i\right)  \\
             & = \frac{1}{n^2}\sum_{i=1}^{n}\mathrm{Var}(X_i)
            = \frac{1}{n^2}\cdot n\sigma^2
            = \frac{\sigma^2}{n}.
        \end{split}
    \end{equation}

    接着对\eqref{eq:variance-decomposition}两侧取期望：
    \begin{equation}
        \begin{split}
            \mathbb{E}\left[\sum_{i=1}^{n}{(X_i-\bar{X})}^2\right]
             & = \mathbb{E}\left[\sum_{i=1}^{n}{(X_i-\mu)}^2\right]
            - n\,\mathbb{E}\left[{(\bar{X}-\mu)}^2\right]           \\
             & = \sum_{i=1}^{n}\mathbb{E}[{(X_i-\mu)}^2]
            - n\,\mathrm{Var}(\bar{X})                              \\
             & = n\sigma^2 - n\cdot\frac{\sigma^2}{n}
            = (n-1)\sigma^2.
        \end{split}
    \end{equation}

    于是
    \begin{equation}
        \begin{split}
            \mathbb{E}[S^2]
             & = \mathbb{E}\left[\frac{1}{n-1}\sum_{i=1}^{n}{(X_i-\bar{X})}^2\right]   \\
             & = \frac{1}{n-1}\,\mathbb{E}\left[\sum_{i=1}^{n}{(X_i-\bar{X})}^2\right] \\
             & = \frac{1}{n-1}\cdot (n-1)\sigma^2
            = \sigma^2.
        \end{split}
    \end{equation}
    这说明 $S^2$ 是总体方差 $\sigma^2$ 的无偏估计。

    同理可得
    \begin{equation}
        \begin{split}
            \mathbb{E}[\mathcal{T}^2]
             & = \mathbb{E}\left[\frac{1}{n}\sum_{i=1}^{n}{(X_i-\bar{X})}^2\right] \\
             & = \frac{1}{n}\,(n-1)\sigma^2
            = \frac{n-1}{n}\sigma^2.
        \end{split}
    \end{equation}
    因此，分母用 $n$ 时会系统性低估总体方差，偏差因子为 $\frac{n-1}{n}$。

\end{AI}

这个证明的思路非常简单，把不知道的量转化为知道的量，凑出来。

\section{多维随机变量与协方差矩阵}
当随机变量是多维的，我们可以使用协方差矩阵来描述它们之间的关系。
设 $X = (X_1, X_2, \ldots, X_n)$ 是一个 $n$ 维随机变量，
那么它的协方差矩阵 $\Sigma$ 定义如\eqref{eq:covariance-matrix}所示。

\begin{equation}
    \Sigma_{ij} = \mathrm{Cov}(X_i, X_j) = \mathbb{E}[(X_i - \mu_i)(X_j - \mu_j)]
    \label{eq:covariance-matrix}
\end{equation}

其中 $\mu_i = \mathbb{E}[X_i]$ 是第 $i$ 个随机变量的均值。
协方差矩阵 $\Sigma$ 是一个对称的正半定矩阵，
它的对角线元素 $\Sigma_{ii}$ 表示第 $i$ 个随机变量的方差，
而非对角线元素 $\Sigma_{ij}$ 表示第 $i$ 和第 $j$ 个随机变量之间的协方差。
协方差矩阵在多维统计分析、机器学习和数据科学等领域中具有重要的应用，
例如在主成分分析（PCA）中，协方差矩阵的特征值和特征向量用于降维和特征提取。

\end{document}